\section{Limites em espaços métricos}

Como na seção anterior, ao longo desta seção, a menos que seja explicitamente dito o contrário, $(M, d)$ e $(N, d')$ são espaços métricos arbitrários.

Na seção anterior, vimos que a continuidade de uma função é trivialmente verdadeira nos pontos isolados de seu domínio.

Nos pontos não isolados de seu domínio, ou seja, nos pontos de seu domínio que são pontos de acumulação, a continuidade de uma função pode ser ou não ser verdadeira, dependendo da função.

Uma ferramenta muito útil para testar a continuidade de uma função nesses pontos é a noção de limite.
Tal noção também dá informações sobre o comportamento de uma função nas imediações de pontos de acumulação de seu domínio que estejam em um espaço ambiente, mas que não pertençam ao domínio.

Intuitivamente, $L$ é limite de uma função $f$ no ponto $p$ se, para qualquer margem de erro em torno de $L$, existe uma margem de segurança em torno de $p$ tal que, para todo ponto do domínio de $f$ dentro dessa margem de segurança e distinto de $p$, a imagem desse ponto por $f$ está dentro da margem de erro em torno de $L$.
O ponto $p$ é excluído dessa análise, pois a noção de limite pergunta o que ocorre perto do ponto $p$, mas não exatamente no ponto $p$.

\begin{definition}\index{limite}\index{limite!definição}
    Sejam
    \begin{itemize}
        \item $M, N$ espaços métricos,
        \item $A\subseteq M$,
        \item $f: A \to N$,
        \item $p \in M$ um ponto de acumulação de $A$,
        \item $L \in N$.
    \end{itemize}
    Dizemos que $L$ é limite de $f$ em $p$ se, para todo $\epsilon > 0$, existe $\delta > 0$ tal que, para todo $x \in A\setminus\{p\}$,
    \[
        \text{se } d(x, p) < \delta, \text{ então } d'(f(x), L) < \epsilon.
    \]
\end{definition}

Conforme veremos mais adiante, nem sempre uma função possui limites em pontos dados.
Porém, limites, quando existem, são únicos.

A Figura~\ref{fig:limiteUnico} resume a escolha de $\delta$ na prova da Proposição~\ref{prop:limiteUnico}.

\begin{proposition}\index{limite!unicidade}\label{prop:limiteUnico}
    Sejam
    \begin{itemize}
        \item $M, N$ espaços métricos,
        \item $A\subseteq M$,
        \item $f: A \to N$,
        \item $p \in M$ um ponto de acumulação de $A$,
        \item $L_1, L_2 \in N$.
    \end{itemize}

    Se $L_1$ e $L_2$ são limites de $f$ em $p$, então $L_1 = L_2$.
\end{proposition}

\begin{proof}
    Suponha por absurdo que $L_1\neq L_2$.
    Então $d'(L_1, L_2) > 0$.

    Seja $R=\frac{d'(L_1, L_2)}{2} > 0$.
    
    Para $i\in \{1, 2\}$, como $L_i$ é limite de $f$ em $p$, existe $\delta_i > 0$ tal que, para todo $x \in A\setminus\{p\}$, se $d(x, p) < \delta_i$, então $d'(f(x), L_i) < R$.

    Seja $\delta = \min\{\delta_1, \delta_2\}$.
    Como $p$ é ponto de acumulação de $A$, existe $x \in A\setminus\{p\}$ tal que $d(x, p) < \delta$.
    Assim, $d(x, p)<\delta\leq \delta_1, \delta_2$, logo:

    \[
        d'(f(x), L_1) < R \qquad\text{e}\qquad d'(f(x), L_2) < R.
    \]
    Assim, da desigualdade triangular, temos que:
    \[
        d'(L_1, L_2) \leq d'(L_1, f(x)) + d'(f(x), L_2) < R + R = 2R = d'(L_1, L_2).
    \]
    Como $d'(L_1, L_2) < d'(L_1, L_2)$ é impossível, chegamos a uma contradição. Logo, $L_1 = L_2$.
\end{proof}


\subimport{limitesMetricos}{limiteUnico}

Com isso, podemos definir uma notação muito usual para se falar de limites.
\begin{definition}\index{limite!notação}
    Seja $A\subseteq M$ um conjunto, $f: A \to N$ uma função e $p \in M$ um ponto de acumulação de $A$.

    Então, caso exista, o único limite de $f$ em $p$ é denotado por
    
    \[
    \lim_{x\to p} f(x).
    \]
\end{definition}
A variável $x$, na notação acima, é uma variável muda.
Ou seja, ela funciona como parte da notação, e não como a representação de algum objeto específico e independente.
Pensamos em $x$ como um ponto do domínio que se aproxima de $p$, mas a letra $x$ pode ser alterada por qualquer outra letra sem que isso altere o significado da notação.
Dessa forma, quando o limite existe,
\[
    \lim_{x\to p} f(x) = \lim_{\alpha\to p} f(\alpha) = \lim_{u\to p} f(u).
\]

No caso de funções de várias variáveis, é comum usar notações alternativas quando conveniente.
Por exemplo, se $f:\mathbb R^2 \to \mathbb R$ é uma função e $p=(a, b)$, é comum usar a notação
\[
    \lim_{(x, y)\to (a, b)} f(x, y).
\]

Para dimensões maiores, notações análogas são utilizadas.
\smallskip


Observe que continuidade em $p$ só faz sentido quando $p$ pertence ao domínio da função.
Já o limite em $p$ pode fazer sentido mesmo quando $p$ não pertence ao domínio, desde que $p$ seja ponto de acumulação do domínio.

Assim, quando $p\in A$, podemos comparar o limite de $f$ em $p$ com o valor $f(p)$.
Quando $p\notin A$, podemos estudar o limite de $f$ em $p$, mas não faz sentido perguntar se $f$ é contínua em $p$.

A proposição abaixo conecta os conceitos de limite e de continuidade.
\begin{proposition}\index{continuidade!caracterização por limite}\label{prop:continuidadeLimite}
    Sejam:
    \begin{itemize}
        \item $A\subseteq M$ um subespaço,
        \item $f: A \to N$ uma função, e
        \item $p \in A$ um ponto de acumulação de $A$.
    \end{itemize}

    São equivalentes:
    \begin{enumerate}[label=(\roman*)]
        \item $f$ é contínua em $p$;
        \item $f(p)$ é limite de $f$ em $p$.
    \end{enumerate}
\end{proposition}

\begin{proof}
    Suponha que $f$ é contínua em $p$.
    Veremos que $f(p)$ é limite de $f$ em $p$.

    Para tanto, fixe $\epsilon>0$.
    Devemos ver que existe $\delta>0$ tal que, para todo $x \in A\setminus\{p\}$, se $d(x, p) < \delta$, então $d'(f(x), f(p)) < \epsilon$.
    Ora, como $f$ é contínua em $p$, temos que tal $\delta$ existe.
    \smallskip
    
    Reciprocamente, suponha que $f(p)$ é limite de $f$ em $p$.
    Devemos ver que $f$ é contínua em $p$.

    Para isso, fixamos $\epsilon>0$.
    Como $f(p)$ é limite de $f$ em $p$, existe $\delta>0$ tal que, para todo $x \in A\setminus\{p\}$, se $d(x, p) < \delta$, então $d'(f(x), f(p)) < \epsilon$.

    Veremos que tal $\delta$ funciona para a definição de continuidade.

    De fato, devemos ver que se $x \in A$ e $d(x, p) < \delta$, então $d'(f(x), f(p)) < \epsilon$.
    Isso decorre da escolha de $\delta$ caso $x \in A\setminus\{p\}$.
    Já se $x = p$, então $d'(f(x), f(p)) = d'(f(p), f(p)) = 0 < \epsilon$.
\end{proof}

Assim, temos o seguinte valioso teste para se decidir se uma função é contínua em um ponto.

\begin{corollary}
    Sejam:
    \begin{itemize}
        \item $A\subseteq M$ um subespaço,
        \item $f: A \to N$ uma função, e
        \item $p \in A$.
    \end{itemize}
    Então:
    \begin{enumerate}[label=(\roman*)]
        \item se $p$ é um ponto isolado de $A$, então $f$ é contínua em $p$;
        \item se $p$ é um ponto de acumulação de $A$ e $f$ não possui limite em $p$, então $f$ não é contínua em $p$;
        \item se $p$ é um ponto de acumulação de $A$ e $f$ possui limite em $p$ que é diferente de $f(p)$, então $f$ não é contínua em $p$;
        \item se $p$ é um ponto de acumulação de $A$ e $f(p)$ é limite de $f$ em $p$, então $f$ é contínua em $p$.
    \end{enumerate}
\end{corollary}
\begin{proof}
    O item (i) segue do fato de que toda função é contínua nos pontos isolados de seu domínio.

    Suponha agora que $p$ é ponto de acumulação de $A$.
    Pelo resultado anterior, $f$ é contínua em $p$ se, e somente se, $f(p)$ é limite de $f$ em $p$.

    Portanto, se $f$ não possui limite em $p$, então $f$ não é contínua em $p$.
    Se $f$ possui limite em $p$, mas esse limite é diferente de $f(p)$, então, pela unicidade do limite, $f(p)$ não é limite de $f$ em $p$, e, portanto, $f$ não é contínua em $p$.
    Finalmente, se $f(p)$ é limite de $f$ em $p$, então $f$ é contínua em $p$.
\end{proof}


Por fim, registramos uma forma precisa do princípio de que a noção de limite depende apenas do comportamento da função em uma vizinhança perfurada do ponto.

\begin{proposition}[Caráter local do limite]\index{limite!caráter local}\label{prop:localidadeLimite}
    Sejam
    \begin{itemize}
        \item $M, N$ espaços métricos,
        \item $A_1,A_2\subseteq M$,
        \item $f_1:A_1\to N$ e $f_2:A_2\to N$ funções,
        \item $a\in M$ um ponto de acumulação de $A_1$ e de $A_2$, e
        \item $L\in N$.
    \end{itemize}

    Suponha que existe um aberto $U\subseteq M$ tal que $a\in U$,
    \[
        (A_1\cap U)\setminus\{a\}\subseteq (A_2\cap U),
    \]
    e, para todo $x\in (A_1\cap U)\setminus\{a\}$,
    \[
        f_1(x)=f_2(x).
    \]

    Nessas condições, se $L$ é limite de $f_2$ em $a$, então $L$ é limite de $f_1$ em $a$.

    Em particular, se $f_2$ possui limite em $a$, então  $f_1$ possui limite em $a$, e, nesse caso,
    \[
        \lim_{x\to a} f_1(x)=\lim_{x\to a} f_2(x).
    \]
\end{proposition}

\begin{proof}
    Devemos ver que fixado $\epsilon>0$, existe $\delta>0$ tal que, para todo $x \in A_1\setminus\{a\}$, se $d(x, a) < \delta$, então $d'(f_1(x), L) < \epsilon$.

    Fixe $\epsilon>0$.
    Como $L$ é limite de $f_2$ em $a$, existe $\delta_1>0$ tal que, para todo $x\in A_2\setminus\{a\}$, se $d(x,a)<\delta_1$, então
    \[
        d'(f_2(x),L)<\epsilon.
    \]

    Como $U$ é aberto e $a\in U$, existe $\delta_2>0$ tal que $B_M(a,\delta_2)\subseteq U$.
    Seja $\delta=\min\{\delta_1,\delta_2\}$.

    Tome $x\in A_1\setminus\{a\}$ tal que $d(x,a)<\delta$.
    Como $d(x,a)<\delta_2$, temos $x\in U$.
    Logo,
    \[
        x\in (A_1\cap U)\setminus\{a\}\subseteq (A_2\cap U)\setminus\{a\}.
    \]
    Portanto, $x\in A_2\setminus\{a\}$ e $f_1(x)=f_2(x)$.
    Como $d(x,a)<\delta_1$, segue que
    \[
        d'(f_1(x),L)=d'(f_2(x),L)<\epsilon.
    \]

    Assim, $L$ é limite de $f_1$ em $a$.
\end{proof}

Em particular, alterar, remover ou acrescentar o valor de uma função no próprio ponto $a$ não altera seu limite em $a$, desde que o limite continue fazendo sentido.

O exemplo abaixo mostra como limites podem ser usados para provar mais facilmente que certas funções não são contínuas em alguns de seus pontos.
Perceba que não foi necessário recorrer à definição que envolve épsilons e deltas.

\begin{example}\label{ex:limiteNaoImporta}
    Seja $f:\mathbb R^2\rightarrow \mathbb R$ dada por
    \[
        f(x, y) = \begin{cases}
            1, & \text{se } (x, y) \neq (0, 0),\\
            0, & \text{se } (x, y) = (0, 0).
        \end{cases}
    \]

    A Figura~\ref{fig:limiteNaoImporta} ilustra o gráfico de $f$.

    \subimport{limitesMetricos}{limiteNaoImporta}

    Seja $1_{\mathbb R^2}:\mathbb R^2\to\mathbb R$ a função constante identicamente igual a $1$.

    A função $f$ não é contínua em $(0, 0)$.
    De fato, primeiro note que $(0, 0)$ é ponto de acumulação de $\mathbb R^2$.
    Como $f$ e $1_{\mathbb R^2}$ coincidem em $\mathbb R^2\setminus\{(0,0)\}$, a Proposição~\ref{prop:localidadeLimite}, a Proposição~\ref{prop:continuidadeLimite} e a continuidade das funções constantes dão
    \[
        \lim_{(x, y)\to (0, 0)} f(x, y)
        =
        \lim_{(x, y)\to (0, 0)} 1_{\mathbb R^2}(x, y)
        =
        1.
    \]

    Assim, $\lim_{(x, y)\to (0, 0)} f(x, y) = 1 \neq f(0, 0)$, e, portanto, $f$ não é contínua em $(0, 0)$.
\end{example}
\subsection*{Exercícios}
\begin{exercise}
    Seja $f:M\rightarrow N$ uma função entre espaços métricos, seja $A\subseteq M$ um subespaço, seja $a\in A$ um ponto de acumulação de $A$ e seja $L\in N$.

    Mostre que $L$ é limite de $f|_A$ em $a$ se, e somente se, para todo $\epsilon>0$ existe $\delta>0$ tal que
    \[
        (f|_A)\left[B_A(a,\delta)\setminus\{a\}\right]\subseteq B_N(L,\epsilon).
    \]
    
    Interprete geometricamente este resultado.
\end{exercise}